\mxcorrections{8}\label{C08-CORR}
\textbf{Usage professeur.} Les réponses qui suivent concernent les tâches non résolues dans le cours. Faire d’abord expliciter la démarche ; fournir une aide graduée avant la solution. Les durées sont des estimations à étalonner en classe.
\subsection*{Diagnostic et remédiation}
\begin{enumerate}
\item $4$. \textit{Action :} Additionner puis diviser par le nombre de valeurs.
\item Oui : nombres positifs de somme 1. \textit{Action :} Mettre au même dénominateur.
\item $9$. \textit{Action :} Garder les parenthèses avant le carré.
\item $\num{0,4}$. \textit{Action :} Utiliser le complément.
\item $5$. \textit{Action :} Écrire gain net = reçu moins payé.
\item $\num{0,5}$. \textit{Action :} Vérifier le carré de la réponse.
\end{enumerate}
\subsection*{Activité de recherche}
Gains $\pt{-1;0;1}$ et $\pt{-3;0;3}$ : espérance 0 dans les deux cas, variances $\dfrac{2}{3}$ et 6. Le second est trois fois plus dispersé en écart type.
\subsection*{Applications immédiates}
\begin{enumerate}
\item Valeurs 0,1,2 ; probabilités $\dfrac{1}{4}$,$\dfrac{1}{2}$,$\dfrac{1}{4}$. $P(X\leq1)=\dfrac{3}{4}$.
\item $E(G)=\dfrac{10}{5}=2$. Pour une mise m, $E(G-m)=2-m$ ; mise équitable 2.
\item $E(X)=0$, $E(X^2)=1$, variance 1 et écart type 1.
\item $E(Y)=-1$, $V(Y)=16$, $\sigma(Y)=4$.
\item $m=\dfrac{6}{5}=\num{1,2}$ et $E(X)=1$ : l’écart vaut \num{0,2} pour cet échantillon.
\end{enumerate}
\subsection*{Exercices : démarche et résultats}
\paragraph{Exercice 8.1 -- Une loi à construire.} Valeurs 0,1,2 ; chacune obtenue par deux faces, probabilités $\dfrac{1}{3}$. $P(X<2)=\dfrac{2}{3}$. Attention au reste 0 pour 3 et 6.
\paragraph{Exercice 8.2 -- Deux calculs de variance.} $E(X)=-\dfrac{1}{2}+1=\dfrac{1}{2}$. $E(X^2)=1+4=5$ ; $V=5-\dfrac{1}{4}=\dfrac{19}{4}$, $\sigma =\dfrac{\sqrt {19}}{2}$. Direct : $(\dfrac{25}{4})(\dfrac{1}{4})+(\dfrac{1}{4})(\dfrac{1}{2})+(\dfrac{49}{4})(\dfrac{1}{4})=\dfrac{76}{16}=\dfrac{19}{4}$.
\paragraph{Exercice 8.3 -- Un jeu équitable.} $E(G)=\dfrac{15}{5}=3$, donc m=3. $E(G^2)=\dfrac{125}{5}=25$ ; $V(G)=25-9=16$, écart type 4, inchangé par la mise.
\paragraph{Exercice 8.4 -- Même moyenne, autre risque.} $E(X)=E(Y)=1$ ; $V(X)=0$ ; $E(Y^2)=10$, $V(Y)=9$. X est certain, Y dispersé. On ne peut pas classer les préférences d’un joueur uniquement à partir de l’espérance.
\paragraph{Exercice 8.5 -- Linéarité sans indépendance.} $E(X)=\dfrac{1+2+3+4+5+6}{6}=\dfrac72$ ; $E(Y)=7-E(X)=\dfrac72$. Comme $X+Y=7$, son espérance vaut 7, somme des deux espérances. Les événements $\{X=1\}$ et $\{Y=6\}$ sont identiques : leur intersection a probabilité $\dfrac16$, tandis que le produit de leurs probabilités vaut $\dfrac1{36}$. Les variables ne sont pas indépendantes. La linéarité reste valable sans cette hypothèse.
\paragraph{Exercice 8.6 -- Une perte minimale.} Écrire $X-a=(X-\mu)+(\mu-a)$, développer puis utiliser $E(X-\mu)=0$. Le carré $(a-\mu)^2$ est minimal, nul seulement pour $a=\mu$. La moyenne minimise l’écart quadratique moyen.
\paragraph{Exercice 8.7 -- Fréquences de lettres.} Passer le texte en minuscules, garder les caractères pour lesquels \code{isalpha()} est vrai, puis compter les « a ». Diviser ce nombre par le nombre de lettres si celui-ci est non nul. Pour un texte sans lettre, renvoyer \code{None}, car la fréquence n’est pas définie. Pour \code{Aba !}, deux « a » sur trois lettres donnent $\dfrac23$ ; les deux autres tests donnent \code{None}. Le programme livré \code{frequence_a} suit cette convention.
\paragraph{Exercice 8.8 -- Reprise différée.} Loi : $0,1,2$ de probabilités $\dfrac{1}{4},\dfrac{1}{2},\dfrac{1}{4}$. Espérance 1 ; moment d’ordre 2 : $\dfrac{1}{2}+1=\dfrac{3}{2}$, variance $\dfrac{1}{2}$. On compte explicitement les quatre chemins.
\paragraph{CQFD, exercice 28 p. 306, questions toutes.} $E(X)=(-10)\times\num{0,35}+(-5)\times\num{0,15}+0\times\num{0,4}+20\times\num{0,1}=-\num{2,25}=-\dfrac94$. Une espérance négative est compatible avec une valeur positive possible.
\paragraph{CQFD, exercice 31 p. 306, questions toutes ; utiliser la loi de l’exercice 28.} $E(X^2)=100\times\num{0,35}+25\times\num{0,15}+400\times\num{0,1}=\num{78,75}$. Ainsi $V(X)=\num{78,75}-(-\num{2,25})^2=\num{73,6875}=\dfrac{1179}{16}$, puis $\sigma(X)=\dfrac{\sqrt{1179}}4\approx\num{8,6}$. Ne pas arrondir l’espérance avant le calcul de variance.
\subsection*{QCM : réponses et raisons}
\textbf{QCM 1.}
\begin{enumerate}
\item B. Partition de l’univers.
\item A. $0+\dfrac{4}{4}$.
\item D. König–Huygens.
\item C. Racine positive.
\item B. Linéarité.
\end{enumerate}
\textbf{QCM 2.}
\begin{enumerate}
\item C. Translation sans dispersion nouvelle.
\item A. Valeur absolue du facteur.
\item D. Une variable peut être négative ou avoir une moyenne non atteinte.
\item B. Univers fini commun suffisant.
\item C. $V(X)=E(X^2)-E(X)^2=5-4=1$.
\end{enumerate}
\subsection*{Vrai ou faux}
\begin{enumerate}
\item Faux : c’est une moyenne de carrés à poids positifs.
\item Faux : constante 1 et variable équiprobable $-2,4$ ont la même moyenne.
\item Vrai : les écarts à la nouvelle moyenne restent les mêmes.
\end{enumerate}
\paragraph{Différenciation.} Pour un blocage technique, proposer la première étape seulement (domaine, facteur ou univers de référence), puis retirer l’aide. Pour un élève autonome, demander une seconde méthode et le traitement d’un cas limite. À l’oral, autoriser une reformulation initiale dans la langue maîtrisée, puis exiger une phrase mathématique précise en français.
